\clearpage\ssr{ВВЕДЕНИЕ}

Задача коммивояжера (Traveling Salesman Problem, TSP) — это классическая задача комбинаторной оптимизации. Её цель — найти маршрут, проходящий через заданное множество городов ровно один раз и минимизирующий общее расстояние или стоимость пути${[1]}$.

Пусть дано множество городов $V = \{v_1, v_2, \dots, v_n\}$ и матрица расстояний $D = [d_{ij}]$, где $d_{ij}$ — расстояние между городами $v_i$ и $v_j$. Требуется найти перестановку $P = (p_1, p_2, \dots, p_n)$, минимизирующую суммарное расстояние:
\[
\min \sum_{k=1}^{n-1} d_{p_k p_{k+1}},
\]
где каждый город посещается ровно один раз. Для замкнутого маршрута дополнительно учитывается возврат в начальный город: $d_{p_n p_1}$.

\textbf{Цель лабораторной работы} --- выполнить сравнительный анализ метода
полного перебора решения задачи коммивояжера с методом на базе муравьиного алгоритма. Для достижения
поставленной цели необходимо выполнить следующие задачи:

\begin{itemize}
	\item[---] описать стандартный метод решения задачи коммивояжера и метод на базе муравьиного алгоритма;
	\item[---] описать модель вычислений;
    \item[---] выполнить оценку трудоемкости разрабатываемых алгоритмов;
    \item[---] реализовать разработанные алгоритмы;
    \item[---] провести сравнительный анализ двух рассмотренных методов решения задачи коммивояжера;
    \item[---] выполнить параметризацию метода, основанного на муравьином алгоритме по трём его параметрам.
\end{itemize}